\documentclass[aspectratio=169,11pt]{beamer}
\usepackage{../../../shared/theme/esmad_beamer_theme}
\usepackage{amsmath,amssymb}
\usepackage{booktabs}

\title[Interference and Spillovers]{Interference, Spillovers, and Network Causal Inference}
\subtitle{When One Unit's Treatment Affects Another Unit's Outcome}
\date{\today}

\begin{document}

\begin{frame}
\titlepage
\end{frame}

\begin{frame}{Learning Goals}
\begin{itemize}
\item explain why SUTVA and no-interference assumptions matter;
\item define causal effects when outcomes depend on treatment assignments of others;
\item distinguish direct, spillover, total, and overall effects;
\item understand exposure mappings and partial interference;
\item connect cluster and two-stage randomization to interference;
\item recognize design and estimation failures caused by contamination and network spillovers.
\end{itemize}
\end{frame}

\begin{frame}{Outline}
\tableofcontents
\end{frame}

\section{Why Interference Matters}

\begin{frame}{Classical Potential Outcomes}
Without interference, unit \(i\) has potential outcomes
\[
Y_i(1),\qquad Y_i(0).
\]

The treatment of other units does not enter unit \(i\)'s outcome.
\end{frame}

\begin{frame}{Treatment Assignment Vectors}
With interference,
\[
Y_i(\mathbf D)
\]
may depend on the entire treatment vector
\[
\mathbf D=(D_1,\ldots,D_n).
\]

\begin{alertblock}{Consequence}
There are potentially many more than two potential outcomes per unit.
\end{alertblock}
\end{frame}

\begin{frame}{Examples of Interference}
\begin{itemize}
\item vaccination protects untreated contacts;
\item pricing in one region shifts demand to neighboring regions;
\item promotions create word-of-mouth effects;
\item agricultural treatment affects nearby plots;
\item information campaigns diffuse through social networks.
\end{itemize}
\end{frame}

\section{SUTVA and No Interference}

\begin{frame}{Stable Unit Treatment Value Assumption}
SUTVA is often summarized as:
\begin{itemize}
\item no multiple versions of treatment;
\item no interference between units.
\end{itemize}

Violations of either part complicate the mapping from treatment labels to potential outcomes.
\end{frame}

\begin{frame}{No-Interference Assumption}
The classical condition is
\[
Y_i(\mathbf D)=Y_i(D_i).
\]

This says assignments received by other units do not matter for unit \(i\)'s outcome.
\end{frame}

\section{Exposure Mappings}

\begin{frame}{Reducing the Assignment Space}
Instead of indexing by the full vector \(\mathbf D\), define an exposure mapping
\[
E_i = h_i(\mathbf D, G),
\]
where \(G\) represents network or group structure.

Then write
\[
Y_i(D_i,E_i).
\]
\end{frame}

\begin{frame}{Examples of Exposure}
Exposure might be:
\begin{itemize}
\item number of treated neighbors;
\item fraction of treated peers;
\item whether at least one neighbor is treated;
\item weighted exposure by network distance;
\item cluster-level treatment saturation.
\end{itemize}
\end{frame}

\begin{alertblock}{Modeling choice}
The exposure mapping is itself an assumption about how interference operates.
\end{alertblock}

\section{Direct and Spillover Effects}

\begin{frame}{Direct Effect}
For fixed peer exposure \(e\), define
\[
DE(e)
=
\mathbb E\left[
Y_i(1,e)-Y_i(0,e)
\right].
\]

This compares own treatment while holding exposure from others fixed.
\end{frame}

\begin{frame}{Spillover Effect}
For an untreated unit,
\[
SE(e_1,e_0)
=
\mathbb E\left[
Y_i(0,e_1)-Y_i(0,e_0)
\right].
\]

This captures how others' treatment affects an untreated unit.
\end{frame}

\begin{frame}{Total and Overall Effects}
A total effect can compare
\[
Y_i(1,e_1)-Y_i(0,e_0),
\]
changing both own treatment and exposure.

An overall effect compares average outcomes under two treatment-allocation regimes.
\end{frame}

\section{Partial Interference}

\begin{frame}{Interference Within Clusters}
Under partial interference:
\begin{itemize}
\item units may affect others within the same cluster;
\item there is no interference across clusters.
\end{itemize}

This makes cluster-level analysis possible while relaxing individual no-interference.
\end{frame}

\begin{frame}{Examples of Clusters}
\begin{itemize}
\item households;
\item schools;
\item villages;
\item hospitals;
\item geographic markets.
\end{itemize}
\end{frame}

\section{Two-Stage Randomization}

\begin{frame}{Randomize Saturation, Then Treatment}
A two-stage design can:
\begin{enumerate}
\item randomize clusters to treatment saturation levels;
\item randomize individuals within clusters to treatment.
\end{enumerate}

This creates variation in both own treatment and peer exposure.
\end{frame}

\begin{frame}{Why Two-Stage Designs Help}
They can identify:
\begin{itemize}
\item direct effects at a given saturation;
\item spillover effects among untreated units;
\item total effects;
\item overall effects of allocation regimes.
\end{itemize}
\end{frame}

\section{Cluster Randomization}

\begin{frame}{Why Cluster-Level Assignment?}
Cluster randomization can reduce contamination when interference is expected within groups.

But it changes:
\begin{itemize}
\item the unit of assignment;
\item the effective sample size;
\item the relevant variance estimator.
\end{itemize}
\end{frame}

\begin{frame}{Effective Sample Size}
With clustered assignment, the number of independent assignment units may be the number of clusters, not the number of individuals.

Few clusters can therefore produce weak inference even with many observations.
\end{frame}

\section{Network Interference}

\begin{frame}{Known Network Structure}
Let \(G_{ij}=1\) indicate that units \(i\) and \(j\) are connected.

Exposure can be defined as
\[
E_i
=
\sum_{j\neq i}G_{ij}D_j.
\]

Weighted or normalized alternatives are also possible.
\end{frame}

\begin{frame}{Distance-Based Interference}
Spillovers may decay with network or geographic distance.

A useful question is not simply whether interference exists, but how far it propagates.
\end{frame}

\begin{frame}{Network Position}
Treatment effects may vary with:
\begin{itemize}
\item degree;
\item centrality;
\item community membership;
\item local treatment density.
\end{itemize}

This mixes interference with heterogeneous treatment effects.
\end{frame}

\section{Contamination and Displacement}

\begin{frame}{Contamination}
Control units can indirectly receive treatment exposure.

Examples:
\begin{itemize}
\item untreated users see treated users' messages;
\item control stores face promotional spillovers;
\item untreated communities interact with treated communities.
\end{itemize}

Naive treated-versus-control comparisons can then understate direct effects.
\end{frame}

\begin{frame}{Displacement}
Treatment can improve outcomes in one unit by worsening outcomes elsewhere.

Examples:
\begin{itemize}
\item crime displaced geographically;
\item sales shifted across stores;
\item patients shifted across hospitals.
\end{itemize}

A local positive effect may coexist with a zero or negative system-wide effect.
\end{frame}

\section{Identification}

\begin{frame}{The Estimand Must Match the Exposure Model}
Credible analysis requires:
\begin{itemize}
\item a treatment assignment mechanism;
\item an interference structure;
\item an exposure mapping;
\item a clearly defined causal contrast.
\end{itemize}

Without these, “the treatment effect” is ambiguous.
\end{frame}

\begin{frame}{Positivity Under Interference}
We need support for relevant combinations of own treatment and exposure.

For example, if highly connected units are always heavily exposed, low-exposure counterfactuals for them may be unsupported.
\end{frame}

\section{Diagnostics}

\begin{frame}{Check the Network / Cluster Assumptions}
Useful diagnostics include:
\begin{itemize}
\item cross-cluster interaction rates;
\item geographic spillover plots;
\item exposure distributions;
\item treatment saturation by cluster;
\item sensitivity to alternative exposure mappings.
\end{itemize}
\end{frame}

\begin{frame}{Alternative Exposure Mappings}
Compare results under:
\begin{itemize}
\item treated-neighbor count;
\item treated-neighbor share;
\item binary any-exposure indicator;
\item distance-weighted exposure.
\end{itemize}

Large changes imply strong dependence on the assumed interference model.
\end{frame}

\section{Applied Reasoning}

\begin{frame}{Example: Regional Pricing Intervention}
Suppose treatment changes prices in selected regions.

Possible spillovers:
\begin{itemize}
\item customers shop across borders;
\item competitors respond in neighboring regions;
\item promotions diffuse nationally;
\item inventory reallocates across regions.
\end{itemize}

A standard regional DiD may then violate no-interference.
\end{frame}

\begin{frame}{A Credible Interference Workflow}
\begin{enumerate}
\item Define the unit and assignment mechanism.
\item State whether no-interference is plausible.
\item Map likely spillover channels.
\item Define the exposure mapping.
\item Choose direct, spillover, total, or regime-level estimands.
\item Match inference to cluster/network dependence.
\item test sensitivity to alternative exposure definitions.
\item report system-wide as well as local effects when displacement is possible.
\end{enumerate}
\end{frame}

\begin{frame}{Common Failure Modes}
\begin{itemize}
\item assuming SUTVA without substantive justification;
\item defining peer exposure after seeing outcomes;
\item ignoring cross-cluster interaction;
\item treating contaminated controls as unexposed;
\item reporting only direct effects when spillovers are policy-relevant;
\item using individual-level standard errors under cluster assignment.
\end{itemize}
\end{frame}

\begin{frame}{Synthesis: Treatment Can Propagate}
\begin{itemize}
\item No-interference is an identifying assumption, not a default truth.
\item With interference, potential outcomes depend on assignment vectors or exposure mappings.
\item Direct, spillover, total, and overall effects answer different questions.
\item Partial interference and cluster designs provide tractable structure.
\item Network information can encode who is exposed to whom.
\item Contamination and displacement can reverse policy interpretation.
\item The estimand, exposure mapping, assignment mechanism, and dependence structure must be aligned.
\end{itemize}
\end{frame}

\begin{frame}{Selected References}
\small
\begin{itemize}
\item Hudgens, M. G., and Halloran, M. E. (2008). Toward Causal Inference with Interference. \textit{JASA}, 103(482), 832--842.
\item Aronow, P. M., and Samii, C. (2017). Estimating Average Causal Effects Under General Interference. \textit{Annals of Applied Statistics}, 11(4), 1912--1947.
\item Tchetgen Tchetgen, E. J., and VanderWeele, T. J. (2012). On Causal Inference in the Presence of Interference. \textit{Statistical Methods in Medical Research}, 21(1), 55--75.
\item Baird, S. et al. (2018). Optimal Design of Experiments in the Presence of Interference. \textit{Review of Economics and Statistics}, 100(5), 844--860.
\item Athey, S., Eckles, D., and Imbens, G. W. (2018). Exact P-values for Network Interference. \textit{JASA}, 113(521), 230--240.
\end{itemize}
\end{frame}

\contactslide

\end{document}
